%% ========================================================================
%% Appendix C: Glossary
%% ========================================================================

\chapter{Glossary}
\label{app:glossary}

Lampiran ini berisi definisi istilah-istilah teknis yang digunakan dalam buku ini.
Istilah utama dipertahankan dalam bahasa aslinya dan ditulis miring saat muncul
di dalam bab.

%% ------------------------------------------------------------------------
\section*{A}

\begin{description}
    \gitem{ACL (Access Control List)}{Mekanisme kontrol akses yang lebih granular
          dibanding \term{permission} Unix tradisional, sehingga hak akses dapat
          diatur untuk banyak pengguna dan grup.}

    \gitem{Address space}{Ruang alamat yang terlihat oleh sebuah proses. Dalam
          \term{virtual memory}, setiap proses memiliki address space sendiri yang
          terisolasi dari proses lain.}

    \gitem{API (Application Programming Interface)}{Sekumpulan definisi dan protokol
          untuk membangun dan mengintegrasikan software aplikasi.}
\end{description}

%% ------------------------------------------------------------------------
\section*{B}

\begin{description}
    \gitem{Bash (Bourne Again Shell)}{\term{Shell} dan bahasa pemrograman
          \term{command-line} yang menjadi default di banyak distribusi Linux.}

    \gitem{Base image}{Image dasar yang dipakai sebagai fondasi untuk membangun
          image container lain.}

    \gitem{Binary}{Berkas program hasil kompilasi yang siap dijalankan sistem operasi.
          Isinya sudah berbentuk kode mesin, bukan \term{source code}.}

    \gitem{Bootloader}{Program yang memuat sistem operasi ke memori saat komputer
          dinyalakan. Contoh umum: GRUB dan LILO.}
\end{description}

%% ------------------------------------------------------------------------
\section*{C}

\begin{description}
    \gitem{Cache}{Penyimpanan sementara berkecepatan tinggi yang menyimpan data atau
          instruksi yang sering diakses agar proses berikutnya lebih cepat.}

    \gitem{CLI (Command Line Interface)}{Antarmuka berbasis teks untuk berinteraksi
          dengan sistem operasi menggunakan perintah.}

    \gitem{Container}{Lingkungan eksekusi terisolasi di tingkat sistem operasi yang
          berbagi kernel host, tetapi memiliki proses, filesystem, dan konfigurasi
          sendiri.}

    \gitem{Cron}{\term{Daemon} untuk menjadwalkan eksekusi perintah atau skrip pada
          waktu tertentu.}
\end{description}

%% ------------------------------------------------------------------------
\section*{D}

\begin{description}
    \gitem{Daemon}{Program yang berjalan di latar belakang dan biasanya tidak memiliki
          antarmuka langsung dengan pengguna.}

    \gitem{Disk layout}{Susunan partisi, \term{filesystem}, dan area penting lain pada
          sebuah disk atau perangkat penyimpanan.}

    \gitem{Driver}{Komponen software yang memungkinkan sistem operasi berkomunikasi
          dengan perangkat keras tertentu melalui antarmuka yang terstandar.}

    \gitem{Distro (Distribution)}{Varian Linux yang dikemas dengan kernel,
          utilitas, dan software aplikasi tertentu.}

\gitem{Dockerfile}{Berkas resep yang berisi langkah-langkah untuk membangun
          sebuah \term{image} container.}
\end{description}

%% ------------------------------------------------------------------------
\section*{E}

\begin{description}
    \gitem{Endpoint}{Titik akhir komunikasi tempat data dikirim atau diterima.
          Dalam konteks jaringan atau IPC, endpoint sering diwujudkan sebagai
          \term{socket}.}
\end{description}

%% ------------------------------------------------------------------------
\section*{F}

\begin{description}
    \gitem{File}{Unit data bernama yang disimpan di \term{filesystem}. Di Linux,
          file bisa berupa dokumen biasa, konfigurasi, perangkat, atau bentuk khusus lain.}

    \gitem{File descriptor}{Nomor identitas kecil yang dipakai kernel untuk
          mewakili berkas, pipe, socket, atau aliran input-output lain yang sedang
          dibuka proses.}

    \gitem{Filesystem}{Struktur dan aturan yang dipakai sistem operasi untuk
          menyimpan, menamai, dan mengatur berkas di media penyimpanan.}

    \gitem{Function}{Blok perintah atau kode yang diberi nama agar dapat dipanggil
          berulang kali untuk menjalankan tugas tertentu.}

    \gitem{Firmware}{Perangkat lunak tingkat rendah yang tertanam di perangkat keras
          untuk menjalankan fungsi dasar sebelum sistem operasi aktif penuh.}
\end{description}

%% ------------------------------------------------------------------------
\section*{G}

\begin{description}
    \gitem{Group}{Kumpulan akun pengguna yang dipakai Linux untuk mengatur hak akses
          bersama terhadap berkas dan sumber daya sistem.}
\end{description}

%% ------------------------------------------------------------------------
\section*{H}

\begin{description}
    \gitem{Host}{Sistem utama tempat program, \term{virtual machine}, atau
          \term{container} dijalankan. Dalam jaringan, host juga dapat berarti mesin
          yang memiliki identitas alamat tersendiri.}
\end{description}

%% ------------------------------------------------------------------------
\section*{I}

\begin{description}
    \gitem{Image}{Cetakan dasar yang bersifat statis untuk membuat satu atau lebih
          \term{container}. Dalam Docker, image dibangun berlapis dari instruksi di
          \term{Dockerfile}.}

    \gitem{Init system}{Sistem pertama yang berjalan di \term{user space} setelah
          kernel selesai melakukan boot. Tugasnya memulai dan mengelola layanan lain.}

    \gitem{Interrupt}{Sinyal dari perangkat keras atau software yang meminta CPU
          menghentikan sementara tugas saat ini untuk menangani kejadian tertentu.}
\end{description}

%% ------------------------------------------------------------------------
\section*{J}

\begin{description}
    \gitem{Journal}{Sistem pencatatan log terstruktur, biasanya merujuk pada log yang
          dikelola oleh \term{systemd-journald}.}
\end{description}

%% ------------------------------------------------------------------------
\section*{K}

\begin{description}
    \gitem{Kernel}{Inti dari sistem operasi yang mengelola hardware dan menyediakan
          layanan dasar untuk software lain.}

    \gitem{Kernel space}{Wilayah eksekusi dengan hak istimewa penuh tempat kernel dan
          sebagian komponen inti sistem berjalan. Akses ke hardware dikendalikan dari sini.}
\end{description}

%% ------------------------------------------------------------------------
\section*{L}

\begin{description}
    \gitem{Latency}{Waktu tunda antara permintaan dan respons. Dalam sistem operasi,
          latency sering dipakai untuk menilai seberapa cepat sistem bereaksi.}
\end{description}

%% ------------------------------------------------------------------------
\section*{M}

\begin{description}
    \gitem{Mount}{Proses menghubungkan \term{filesystem}, partisi, atau perangkat
          penyimpanan ke satu titik direktori agar isinya bisa diakses dari pohon direktori Linux.}
\end{description}

%% ------------------------------------------------------------------------
\section*{O}

\begin{description}
    \gitem{Override}{Konfigurasi yang menimpa atau mengganti nilai bawaan tanpa
          mengubah berkas asli sumbernya. Dalam \term{systemd}, override lokal biasanya
          disimpan terpisah dari unit bawaan paket.}
\end{description}

%% ------------------------------------------------------------------------
\section*{P}

\begin{description}
    \gitem{Package manager}{Perangkat lunak untuk memasang, memperbarui, menghapus,
          dan melacak dependensi paket aplikasi. Contoh: \cmd{apt} dan \cmd{dpkg}.}

    \gitem{Page fault}{Kejadian ketika proses mengakses halaman memori virtual yang
          belum tersedia di RAM, sehingga kernel harus memetakan atau memuatnya lebih dulu.}

    \gitem{Page table}{Struktur data yang dipakai kernel untuk memetakan alamat
          \term{virtual memory} ke alamat memori fisik.}

    \gitem{Partition layout}{Susunan partisi pada sebuah disk, termasuk urutan,
          ukuran, dan fungsi setiap partisi.}

    \gitem{Paging}{Teknik manajemen memori yang membagi memori virtual dan fisik ke
          dalam blok-blok kecil bernama page agar pemetaan lebih fleksibel.}

    \gitem{Permission}{Aturan hak akses baca, tulis, dan eksekusi pada berkas atau
          direktori di Linux.}

    \gitem{PID (Process ID)}{Nomor identitas unik yang diberikan sistem operasi kepada
          setiap proses yang sedang berjalan.}

    \gitem{Pipeline}{Teknik menghubungkan output satu perintah ke input perintah lain
          menggunakan operator \texttt{|}.}
\end{description}

%% ------------------------------------------------------------------------
\section*{R}

\begin{description}
    \gitem{Redirection}{Teknik mengalihkan input atau output perintah ke berkas,
          perangkat, atau perintah lain menggunakan operator seperti \texttt{>},
          \texttt{>>}, dan \texttt{<}.}

    \gitem{Repository}{Sumber paket software yang menyimpan metadata dan berkas yang
          diperlukan \term{package manager} untuk instalasi dan pembaruan.}

    \gitem{Runlevel}{Konsep lama pada sistem init tradisional untuk menyatakan mode
          operasi sistem, misalnya mode multi-user atau mode pemeliharaan. Dalam
          \term{systemd}, konsep ini banyak digantikan oleh \term{target}.}
\end{description}

%% ------------------------------------------------------------------------
\section*{S}

\begin{description}
    \gitem{Scheduler}{Bagian kernel yang menentukan proses atau thread mana yang
          mendapat giliran memakai CPU, kapan dijalankan, dan berapa lama.}

    \gitem{Service}{Program yang berjalan di latar belakang untuk menyediakan fungsi
          tertentu, misalnya server web, SSH, atau database.}

    \gitem{Shared library}{Pustaka biner yang dapat dipakai bersama oleh banyak program
          saat berjalan, sehingga kode umum tidak perlu disalin ke setiap \term{binary}.}

    \gitem{Shell}{Program yang menyediakan antarmuka bagi pengguna untuk berinteraksi
          dengan sistem operasi, baik melalui perintah teks maupun skrip.}

    \gitem{Signal}{Pesan singkat yang dikirim kernel atau proses lain untuk memberi tahu
          sebuah proses agar berhenti, dijeda, dilanjutkan, atau menangani kejadian tertentu.}

    \gitem{Socket}{Titik komunikasi yang dipakai proses untuk bertukar data, baik
          melalui jaringan maupun komunikasi lokal antarproses.}

    \gitem{Source code}{Kode program yang ditulis manusia dalam bahasa pemrograman dan
          harus diproses lebih dulu sebelum menjadi \term{binary} yang bisa dijalankan.}

    \gitem{Swap}{Ruang di disk yang dipakai sebagai perluasan RAM ketika memori fisik
          tidak cukup untuk menampung seluruh kebutuhan proses aktif.}

    \gitem{System call}{Antarmuka resmi yang dipakai program di \term{user space}
          untuk meminta layanan dari kernel sistem operasi.}

    \gitem{Systemd}{\term{Init system} modern yang digunakan untuk melakukan bootstrap
          \term{user space} dan mengelola layanan sistem.}
\end{description}

%% ------------------------------------------------------------------------
\section*{T}

\begin{description}
    \gitem{Throughput}{Jumlah pekerjaan yang bisa diselesaikan sistem dalam satu satuan
          waktu, misalnya jumlah request per detik atau proses per menit.}
\end{description}

%% ------------------------------------------------------------------------
\section*{U}

\begin{description}
    \gitem{Unit}{Objek konfigurasi yang dikelola \term{systemd}, seperti
          \texttt{.service}, \texttt{.socket}, \texttt{.mount}, dan \texttt{.target}.}

    \gitem{User space}{Wilayah eksekusi tempat aplikasi biasa berjalan dengan hak akses
          terbatas. Program di sini harus memakai \term{system call} untuk meminta layanan kernel.}
\end{description}

%% ------------------------------------------------------------------------
\section*{V}

\begin{description}
    \gitem{Virtual memory}{Lapisan abstraksi memori yang membuat setiap proses seolah-olah
          memiliki ruang alamatnya sendiri, terpisah dari proses lain dan dari memori fisik.}

    \gitem{Virtual machine}{Komputer virtual yang disimulasikan di dalam komputer fisik.
          Ia memiliki sistem operasi, memori, disk, dan perangkat virtualnya sendiri.}

    \gitem{Volume}{Mekanisme penyimpanan data persisten pada container yang dapat
          dipasang dari \term{host} atau dikelola langsung oleh runtime container.}
\end{description}
